\documentclass[11pt]{article}
\usepackage[margin=1in]{geometry}
\usepackage{amsmath,amssymb,hyperref}
\newcommand{\E}{\mathbb E}
\newcommand{\KL}{D_{\mathrm{KL}}}
\title{A Hybrid Discrete--Continuous RePlaid ELBO}
\author{Implementation and derivation notes}
\date{September 29, 2026}
\begin{document}
\maketitle

\paragraph{Scope.} This note derives a proposed combination of CADD's hybrid
reverse kernel and RePlaid's likelihood objective and schedule learning.
It is not a theorem claimed jointly by those papers. The implementation is
\texttt{algos/replaid\_hybrid.py}, with sampler
\texttt{samplers/replaid\_hybrid\_ddpm.py}. The existing RePlaid and
RePlaidKeep implementations are separate baselines.

\section{Forward process and assumptions}
Let $X=(X_1,\ldots,X_n)$ be a discrete sequence and $e_i=E(X_i)$ a
learned, token-local embedding. Embeddings are normalized as in our RePlaid
backbone. Fixed conditioning tokens $C$ are always visible; every formula
below can be read conditional on $C$, with sums over target positions only.
Let $r(t)$ increase from zero to one and $s(t)=\exp(-\gamma(t))$ decrease.
Independently across positions, $M_{t,i}=1$ indicates that a token has been
masked, with probability $r(t)$. Its state is
\[
 Y_{t,i}=\begin{cases}
 X_i,&M_{t,i}=0,\\
 (\mathrm{MASK},Z_{t,i}),&M_{t,i}=1,
 \end{cases}\qquad
 Z_{t,i}=a(t)e_i+b(t)\epsilon_i,
 \quad a^2=\operatorname{sigmoid}(-\gamma),\quad
 b^2=\operatorname{sigmoid}(\gamma).
\]
This specifies a Markov process by using absorbing masking with forward
hazard $r'(t)/(1-r(t))$. At the first masking time, initialize the latent
from the Gaussian marginal at that global time, then continue the usual
variance-preserving Gaussian process. Starting an independent Gaussian
clock at the masking time would define a different process.

The network predicts $\pi_{\theta,i}(v\mid Y_t)$ and uses the same head for
\[
 \widehat e_i(Y_t)=\sum_v\pi_{\theta,i}(v\mid Y_t)E(v).
\]
Visible tokens are copied in the reverse process. The clean-context input
identifies visible positions. Since embeddings have unit norm, its zero
value marks a masked position. Gaussian log-SNR identifies the physical
path coordinate; the model is not conditioned on a separate arbitrary clock.
A contextual encoder $e_i=E_i(X)$ is outside this derivation: revealing
$X_i$ would not determine $e_i$, invalidating the deterministic reveal kernel.

\section{Reverse-step KL and the continuous-time bound}
For $v<t$, conditional on $M_{t,i}=1$, the reverse event probabilities are
\[
 \rho_{\rm reveal}=\frac{r(t)-r(v)}{r(t)},\qquad
 \rho_{\rm stay}=\frac{r(v)}{r(t)}.
\]
Use these same event probabilities in the model. On reveal, predict a
categorical token and set its embedding deterministically. On staying
masked, use the Gaussian bridge variance and replace the true clean
embedding in its mean by $\widehat e_i$. The resulting conditional KL is
\[
 \rho_{\rm reveal}[-\log\pi_{\theta,i}(X_i\mid Y_t)]
 +\rho_{\rm stay}\frac{s(v)-s(t)}{2}\|e_i-\widehat e_i\|^2.
\]
Visible positions have zero KL. Taking $v=t-h$ and $h\downarrow0$ gives
\begin{equation}\label{eq:bound}
 \boxed{\mathcal L=\mathcal L_{\rm prior}+\mathcal L_{\rm rec}
 +\int_0^1\E\sum_i M_{t,i}\left[
 \frac{r'(t)}{r(t)}(-\log\pi_{\theta,i}(X_i\mid Y_t))
 -\frac{s'(t)}2\|e_i-\widehat e_i\|^2\right]dt.}
\end{equation}
This is the continuous-time negative ELBO (an upper bound on expected NLL),
under the standard limiting regularity/integrability assumptions. A finite
step sampler approximates this continuous-time model. It does not make a
finite-step likelihood bound numerically equal to the training integral.

We choose $r(0)=0$: reconstruction at zero is deterministic and
$\mathcal L_{\rm rec}=0$. With $r(1)=1$ the prior is all masked with
independent standard Gaussian latents. For embedding dimension $d$ its
per-position KL is retained exactly:
\[
 \mathcal L_{{\rm prior},i}=\tfrac12\left[
 a_1^2\|e_i\|^2+d(b_1^2-1-\log b_1^2)\right].
\]
Finite terminal SNR does not exactly match the prior; this term accounts
for it and supplies endpoint gradients. No forward-distribution-dependent
``constants'' are dropped while learning embeddings or endpoints.

Setting $r\equiv1$ removes the reveal term and recovers the continuous
Plaid/RePlaid form, with its nonzero endpoint reconstruction restored.
Setting $s\equiv0$ removes the Gaussian denoising term and recovers the
absorbing discrete objective. CADD's practical CE-only surrogate is not
Equation~\eqref{eq:bound}: we keep both weighted terms.

\section{Information identity and the learned clock}
Assume exact posterior categorical predictions and posterior-mean embedding
predictions. Define $\ell^*(t)$ as the expected integrand in
Equation~\eqref{eq:bound}. Then
\begin{equation}\label{eq:info}
 \ell^*(t)=\frac{d}{dt}H(X\mid Y_t)=-\frac{d}{dt}I(X;Y_t).
\end{equation}
For the jump contribution, hold all other observations $B=Y_{t,-i}$ fixed.
Replacing visible $X_i$ by its noisy observation $Z_i$ increases expected
conditional entropy by
\[
 H(X\mid B,Z_i)-H(X\mid B,X_i)=H(X_i\mid B,Z_i).
\]
The equality uses $Z_i\perp X_{-i}\mid X_i,B$, true for token-local
independent Gaussian channels. Multiplying by jump hazard and visible
probability yields $r'(t)\E[H(X_i\mid Y_t)\mid M_{t,i}=1]$.
This equals the masked CE contribution in~\eqref{eq:bound}.
For Gaussian evolution at fixed mask pattern, conditional vector I--MMSE
gives $-s'(t)/2$ times the summed posterior embedding MSE at masked positions.
Adding these first-order contributions proves~\eqref{eq:info}; simultaneous
masking events have second-order probability away from endpoints.

Specify a fixed path $r=R(u),\gamma=G(u)$ and a monotone clock $u=g(t)$.
Let $A_\theta(u)$ be the expected ELBO density per unit $u$. Then
\[
 \ell_{\theta,g}(t)=g'(t)A_\theta(g(t)),\qquad
 \int_0^1\ell_{\theta,g}(t)dt=\int_0^1A_\theta(u)du.
\]
For fixed model, embeddings, and path endpoints, the clock changes only
sampling variance. Where $A_\theta>0$, the ideal mean-flattening clock obeys
\[
 t=\frac{\int_0^{g(t)}A_\theta(u)du}{\int_0^1A_\theta(u)du},
 \qquad \ell_{\theta,g}(t)=\kappa.
\]
Together with Bayes optimality this implies
$I(X;Y_t)=I(X;Y_0)-\kappa t$. For conditional generation replace information
by conditional information given $C$. This does not imply linear
masked-token CE or eliminate sequence dependencies.
Learning two independent path functions changes the path itself and is not
covered by the time-reparameterization invariance.

\section{Implemented estimator and unbiased clock gradient}
We use $R(u)=u$, $G(u)=\gamma_0+(\gamma_1-\gamma_0)u$ and ordered learnable
endpoints, with positive width represented by softplus. Endpoint parameters
and embeddings minimize the full NELBO. There is no extra CE coefficient.

Direct masked estimation contains $M_i\,\mathrm{CE}_i/u$. Its second moment
can diverge near zero when the predictor has nonzero clean-end CE. Instead
choose a target position $J$ uniformly among the $N$ loss positions and
sample the corruption conditional on $M_J=1$. Other masks remain independent
Bernoulli($u$). With $w(u)=\tfrac12(\gamma_1-\gamma_0)e^{-G(u)}$,
\[
 \widehat A(u)=\mathrm{CE}_J+u\,w(u)\|e_J-\widehat e_J\|^2,
 \qquad \E\widehat A(u)=\frac1N A_\theta(u).
\]
This follows from $\E[M_i f(Y)]=u\E[f(Y)\mid M_i=1]$ and uniform target
selection. It remains unbiased for arbitrarily correlated tokens. The initial implementation retained only this selected-target loss. The
corrected implementation Rao--Blackwellizes over the selected target:
given the observed mask with $K$ masked positions, the posterior of $J$
is uniform on those positions. Therefore use
\[
 \widehat A_{\rm RB}(u)=\frac1K\sum_{i:M_i=1}
 [\mathrm{CE}_i+u w(u)\|e_i-\widehat e_i\|^2].
\]
The forced-target proposal has mask density ratio $K/(Nu)$ relative to the
original Bernoulli process. Thus this estimator is unbiased and has no larger
conditional variance than retaining the randomly chosen target. Code assigns
$N/K$ times each masked-position contribution before valid-token reduction.
This uses all predictions from the same forward pass, with no extra network
compute. All-context rows contribute zero. Binary loss
masks are required, and loss-excluded positions are fixed conditioning.

Represent a density $p_\phi(u)$ with $K$ equal-width bins and probabilities
\[
 P_k=(1-\eta)\operatorname{softmax}(h)_k+\eta/K,
 \qquad p_\phi(u)=KP_k \text{ in bin }k.
\]
We use $K=64$, $\eta=0.01$. The clock is $g_\phi=F_\phi^{-1}$; its derivative
is $1/p_\phi(g_\phi(t))$ almost everywhere. Training draws one random point in each equal-probability clock stratum,
randomly permutes those points across examples, and applies the inverse CDF.
Each row retains the correct marginal density; stratification reduces time
sampling noise. The likelihood estimator is
$L=\widehat A(u)/p_\phi(u)$, using detached density in likelihood gradients.
Sampling uses the same inverse CDF to set both masking and Gaussian noise.

For fixed non-clock parameters, the second moment and its gradient are
\[
 J(\phi)=\int\frac{\E[\widehat A(u)^2]}{p_\phi(u)}du,
 \qquad
 \nabla_\phi J=-\E_{u\sim p_\phi}
 [L^2\nabla_\phi\log p_\phi(u)].
\]
The implementation adds a zero-valued surrogate
$S-\operatorname{stopgrad}(S)$, where
$S=-\operatorname{stopgrad}(L^2)\log p_\phi(u)$.
It changes only clock gradients, not reported NELBO values, model gradients,
or endpoint gradients. This avoids incorrect pathwise differentiation
through sampled discrete masks. The prior does not depend on the clock and
is excluded from its variance objective.

\paragraph{Exact versus approximate information linearity.}
The ideal mean-flattening density is proportional to $A_\theta(u)$.
The per-example stochastic second-moment optimum instead satisfies
$p^*(u)\propto\sqrt{\E[\widehat A(u)^2]}$. The implemented clock score minimizes this per-example second moment, not
the full stratified batch variance with its cross-example covariances.
Conditional sampling noise,
finite bins, the uniform floor, and an imperfect denoiser prevent a claim of
exact linear information in this implementation. The mean-level theorem
survives; empirical linearity is a diagnostic, not a guaranteed outcome.

\section{Sampling and verification}
Start from all-masked standard Gaussian latents. At reverse clock times
$t>v$, compute $u_t=g(t),u_v=g(v)$, reveal with probability
$(u_t-u_v)/u_t$, and draw revealed tokens categorically. Remaining continuous
positions use the Gaussian bridge with the same endpoint schedule and soft
embedding prediction. Already visible positions and prompts remain fixed.
The grid ends exactly at zero, where all remaining tokens reveal; no extra
argmax decoder is used. Nonzero sampling epsilon is rejected.

Tests cover the importance integral, the exact expected score-gradient
against analytic differentiation, conditional-estimator equivalence by
mask enumeration, endpoint reveal probabilities, singleton batches,
embedding/head/endpoint/clock gradients, clean prompt preservation,
all-context rows, composition of the matched Sudoku configuration, fresh-process sampler
registration, default public generation endpoints, and exact enumeration
of the Rao--Blackwellized mask estimator.
The matched experiment is \texttt{sudoku\_replaid\_hybrid\_matched}, using
W\&B project \texttt{replaid-sudoku}. Implementation validation is not a
claim of Sudoku performance; that requires a separate training run.

\begin{thebibliography}{9}
\bibitem{cadd} Zheng et al. \emph{Continuously Augmented Discrete Diffusion
model for Categorical Generative Modeling}, equations 4.14--4.17.
\url{https://arxiv.org/html/2510.01329}
\bibitem{replaid} \emph{Continuous Diffusion Scales Competitively with
Discrete Diffusion for Language}, Section 5.2.
\url{https://arxiv.org/html/2605.18530v1}
\bibitem{plaid} Gulrajani and Hashimoto. \emph{Likelihood-Based Diffusion
Language Models}. \url{https://arxiv.org/abs/2305.18619}
\end{thebibliography}
\end{document}
